\documentclass[../general/general]{subfiles}

\begin{document}


\chapter{Лекция №3}
\noindent\textit{А.\,М.~Белемук \hfill 27 февраля 2021}
\vspace{1em}


% ==============================================================================
% РАЗДЕЛ 1
% ==============================================================================
\section{Термодинамические следствия канонического ансамбля для идеального газа}

На предыдущей лекции мы вычислили свободную энергию идеального газа через каноническое распределение:
\begin{equation}
    F = -TN \ln \left( \frac{eV}{N V_Q} \right), \qquad V_Q = \lambda_T^3 = \left( \frac{2\pi\hbar^2}{mT} \right)^{3/2}
\end{equation}
Теперь из дифференциала $dF = -S\,dT - P\,dV + \mu\,dN$ получим все термодинамические характеристики газа. Мы работаем в классическом режиме $v = V/N \gg V_Q$.

% --- НАЧАЛО БЛОКА 1.2 ---
\subsection{Вычисление давления. Уравнение состояния идеального газа}

Мы знаем, что свободная энергия является функцией переменных температуры $T$, объема $V$ и числа частиц $N$: $F = F(T, V, N)$. Запишем дифференциал свободной энергии:
\begin{equation}
    dF = -S dT - P dV + \mu dN
\end{equation}
Отсюда мы сразу считываем все частные производные. Начнем с давления $P$:
\begin{equation}
    P = -\left( \frac{\partial F}{\partial V} \right)_{T,N}
\end{equation}
Начинаем дифференцировать нашу свободную энергию. Зависимость от объема находится под логарифмом. Мы можем представить свободную энергию в виде $-TN \ln V$ плюс некоторая функция, зависящая только от температуры $T$ и числа частиц $N$:
\begin{equation}
    -\left( \frac{\partial F}{\partial V} \right)_{T,N} = \frac{\partial}{\partial V} \Big[ T N \ln V + f(T,N) \Big]
\end{equation}
Дифференцируя по $V$, получаем:
\begin{equation}
    P = T N \frac{1}{V}
\end{equation}
Делаем вывод:
\begin{equation}
    P = \frac{N}{V} T = n T
\end{equation}
где $n = N/V$ --- концентрация частиц. Мы получили уравнение состояния классического идеального газа, то, что и следовало ожидать. Зависимость давления от концентрации и температуры.
% --- КОНЕЦ БЛОКА 1.2 ---

% --- НАЧАЛО БЛОКА 1.3 ---
\subsection{Вычисление энтропии. Формула Саккура-Тетроде}

Давайте теперь посмотрим на энтропию $S$. Из дифференциала $dF$ следует, что:
\begin{equation}
    S = -\left( \frac{\partial F}{\partial T} \right)_{V,N}
\end{equation}
Дифференцируем выражение для свободной энергии по температуре (два раза знак минус дает плюс):
\begin{equation}
    S = \frac{\partial}{\partial T} \left\{ T N \ln \left[ \frac{e V}{N} \left( \frac{mT}{2\pi \hbar^2} \right)^{3/2} \right] \right\}
\end{equation}
Здесь у нас произведение $T$ на логарифм, который сам зависит от $T$. Применяем правило дифференцирования произведения. Первое слагаемое, производная по $T$ от множителя $T$, дает просто единицу, умноженную на логарифм:
\begin{equation}
    N \ln \left( \frac{e V}{N} \frac{1}{V_Q} \right)
\end{equation}
Второе слагаемое — это $TN$, умноженное на производную по температуре от логарифма. Зависимость от температуры сидит в множителе $T^{3/2}$. Поэтому мы пишем:
\begin{equation}
    + T N \frac{\partial}{\partial T} \ln \left( T^{3/2} \right) + \text{производная от функции }(V,N)
\end{equation}
Производная от $\ln(T^{3/2})$ дает $\frac{3}{2} \frac{1}{T}$. Умножая это на $T N$, получаем:
\begin{equation}
    T N \left( \frac{3}{2} \frac{1}{T} \right) = \frac{3}{2} N
\end{equation}
В результате энтропия получается равной:
\begin{equation}
    S = N \ln \left( \frac{e V}{N} \frac{1}{V_Q} \right) + \frac{3}{2} N
\end{equation}
Давайте преобразуем это выражение. Мы помним, что $\ln e = 1$. Можем внести $\frac{3}{2}N$ под скобки. Получается:
\begin{equation}
    S = N \left[ \ln \left( \frac{V}{N} \frac{1}{V_Q} \right) + 1 + \frac{3}{2} \right] = N \left[ \ln \left( \frac{v}{V_Q} \right) + \frac{5}{2} \right]
\end{equation}
где $v = V/N$ --- удельный объем. 

Это точное выражение для энтропии идеального газа как функции от $T, V, N$. Называется это выражение \textbf{формула Саккура-Тетроде} (Sackur-Tetrode). Она была установлена до создания квантовой механики, в 1912 году, Сакуром и, отдельно, Тетроде. Вы видите, что энтропия идеального газа определена нами не только с точностью до какой-то константы $S_0$, но и абсолютной величиной. Это очень важный результат.
% --- КОНЕЦ БЛОКА 1.3 ---

% --- НАЧАЛО БЛОКА 1.4 ---
\subsection{Химический потенциал классического идеального газа}

Что нам остается? Давайте посчитаем химический потенциал $\mu$. Я вижу из дифференциала $dF$, что:
\begin{equation}
    \mu = \left( \frac{\partial F}{\partial N} \right)_{T,V}
\end{equation}
Начинаю дифференцировать по $N$:
\begin{equation}
    \mu = \frac{\partial}{\partial N} \left[ -T N \ln \left( \frac{e V}{N} \frac{1}{V_Q} \right) \right]
\end{equation}
Опять дифференцируем произведение. Сначала берем производную по множителю $N$:
\begin{equation}
    -T \ln \left( \frac{e V}{N} \frac{1}{V_Q} \right)
\end{equation}
Затем берем производную по логарифму:
\begin{equation}
    -T N \frac{\partial}{\partial N} \ln \left( \frac{1}{N} \right)
\end{equation}
Производная $\frac{\partial}{\partial N} (-\ln N) = -\frac{1}{N}$. Умножая на $-T N$, получаем просто $+T$.
В итоге:
\begin{equation}
    \mu = -T \ln \left( \frac{e V}{N} \frac{1}{V_Q} \right) + T = -T \left[ \ln\left(\frac{V}{N V_Q}\right) + \ln e - 1 \right] = -T \ln \left( \frac{v}{V_Q} \right)
\end{equation}
Вот такой ответ. То есть мы получаем зависимость химического потенциала идеального газа как функции температуры $T$, объема $V$ и числа частиц $N$:
\begin{equation}
    \mu(T, V, N) = -T \ln \left( \frac{v}{V_Q} \right)
\end{equation}
Это важнейший результат: \textbf{химический потенциал идеального газа}.

Давайте посмотрим на его свойства. Во-первых, мы видим, что химический потенциал классического идеального газа отрицателен: $\mu < 0$. Почему? Так как под логарифмом стоит число, много большее единицы: $v \gg V_Q$ (условие применимости классического режима). 
Во-вторых, мы видим, что с ростом температуры $T$ химпотенциал убывает, и убывает сильно. 

Если нарисовать график зависимости $\mu(T)$:
\begin{center}
    \includegraphics[width=0.55\textwidth]{mu(T)}
\end{center}
Вот она себя ведет так: если провести прямую линию $-T$, то мой химпотенциал сильно от нее отклоняется вниз с ростом температуры. И всё это справедливо, когда $v \gg V_Q$. На языке температур это означает, что $T \gg T_0$, где $T_0$ — температура вырождения. Мы с вами будем оценивать температуру вырождения отдельно для Ферми- и Бозе-газов позже.

Еще одна удобная запись для химпотенциала. Так как концентрация $n = 1/v$, можно переписать:
\begin{equation}
    \mu = T \ln (n V_Q)
\end{equation}
Это тоже очень полезная формула, запомним её.
% --- КОНЕЦ БЛОКА 1.4 ---

% --- НАЧАЛО БЛОКА 1.5 ---
\subsection{Внутренняя энергия и связь со статсуммой}

Осталось найти внутреннюю энергию $E$. Как сразу энергию можно посчитать, зная статсумму, минуя вычисления энтропии и промежуточных шагов? Мы знаем, что:
\begin{equation}
    E = F + TS
\end{equation}
Подставим сюда выражения для $F$ и $S$:
\begin{equation}
    E = F + T \left( -\frac{\partial F}{\partial T} \right)
\end{equation}
Заменим $F$ на $-T \ln Z$:
\begin{equation}
    E = -T \ln Z + T \frac{\partial}{\partial T} (T \ln Z)
\end{equation}
Раскроем производную:
\begin{equation}
    E = -T \ln Z + T \left( \ln Z + T \frac{\partial}{\partial T} \ln Z \right)
\end{equation}
Видите, слагаемые $-T \ln Z$ и $+T \ln Z$ просто сокращаются, и остается:
\begin{equation}
    E = T^2 \frac{\partial}{\partial T} \ln Z
\end{equation}
Очень удобная формула! Наряду с тем, что $F = -T \ln Z$, мы получаем еще одну очень удобную формулу для энергии:
\begin{equation}
    E = T^2 \frac{\partial}{\partial T} \ln Z
\end{equation}
Эта операция $T^2 \frac{\partial}{\partial T}$ еще раз нам встретится, когда мы будем вычислять флуктуации, так что мы возьмем ее на вооружение. 

Давайте ее применим благополучно к нашему выражению. Мы знаем, что статсумма $Z \propto T^{3N/2}$ (остальные множители от температуры не зависят). 
Берем логарифм:
\begin{equation}
    \ln Z = \frac{3N}{2} \ln T + \text{const по } T
\end{equation}
Подставляем в формулу:
\begin{equation}
    E = T^2 \frac{\partial}{\partial T} \left( \frac{3N}{2} \ln T \right) = T^2 \cdot \frac{3N}{2} \frac{1}{T} = \frac{3}{2} N T
\end{equation}
Знакомый результат? Мы получили, что энергия идеального газа есть $\frac{3}{2}NT$. То, что вы уже знали из курса общей физики. Вот так элегантно через статсумму можно получать буквочку $E$.

Ну что, вот мы с каноническим распределением в классической статистической физике познакомились, и теперь приступаем к теме следующей нашей лекции — равновесной матрице плотности.
% --- КОНЕЦ БЛОКА 1.5 ---

% ==============================================================================
% РАЗДЕЛ 2
% ==============================================================================
\section{Аппарат квантовой статистической механики: чистые состояния}

Ну что, вот мы с каноническим распределением в классической статистической физике познакомились, а теперь приступаем к теме следующей нашей лекции номер три. Тема нашей сегодняшней лекции — \textbf{Равновесная матрица плотности}. 

Надо понять, о чем речь, что это такое и для чего ее вводят. Мы приступаем к изучению статистической механики квантовых систем, которая наиболее нам интересна. Именно в квантовых системах газы становятся вырожденными, и чтобы их описывать, вводят новый аппарат. Первый пункт нашего рассмотрения так и будет называться: аппарат квантовой статистической механики. И таким важным аппаратом как раз будет являться матрица плотности. Давайте к ней подбираться постепенно.

% --- НАЧАЛО БЛОКА 2.1 ---
\subsection{Переход от классического фазового пространства к квантовому описанию}

К чему мы привыкли? Была классическая механика. Где мы сразу стали развивать классическую статистическую механику? Мы вводили с вами фазовое пространство импульсов и координат, $6N$-мерное. В этом пространстве у нас двигались изображающие точки.

Теперь переходим к квантовой механике. В квантовой механике импульс $i$-той частицы $\vec{p}_i$ и координата $i$-той частицы $\vec{r}_i$ одновременно неизмеримы. Вы не можете их одновременно измерить и сказать, что система находится в какой-то конкретной точке фазового пространства. Нет одновременно ни точной координаты, ни точного импульса. Можно только в квазиклассическом приближении как-то вводить подобное описание. Что же приходит на смену фазовому пространству? 
% --- КОНЕЦ БЛОКА 2.1 ---

% --- НАЧАЛО БЛОКА 2.2 ---
\subsection{Чистое состояние и чистый ансамбль}

Как описывать состояния квантовой системы? 
%[МЕСТО ДЛЯ ГРАФИКА: Таймкод ~ 25:30. Изображение квантовой системы: лектор рисует прямоугольный ящик (систему), внутри которого схематично нарисованы волны, символизирующие волновую функцию].

Первое и самое простое, к чему мы привыкли — это \textbf{чистое состояние}. 
Что такое чистое состояние? Если состояние системы полностью описывается волновой функцией, то такое состояние называется чистым. Давайте запишем волновую функцию в координатном представлении — это будет функция от всех координат всех частиц:
\begin{equation}
    \Psi(\vec{r}_1, \vec{r}_2, \dots, \vec{r}_N)
\end{equation}
Для сокращенности мы будем обозначать ее просто $\Psi(x)$, где под $x$ дальше будет подразумеваться вся совокупность координат $N$ частиц: $x = \{\vec{r}_1, \vec{r}_2, \dots, \vec{r}_N\}$.

Итак, если состояние системы описывается волновой функцией $\Psi(x)$, то такое состояние называется чистым. Это то, с чем вы работали, когда изучали квантовую механику: вы писали для гамильтониана какую-то волновую функцию. И вы знаете, как эта волновая функция эволюционирует во времени. Вы всё знаете про систему:
\begin{equation}
    \Psi(x, t) = e^{-\frac{i}{\hbar}\hat{H}t} \Psi(x)
\end{equation}
Вы действуете оператором эволюции на начальное состояние $\Psi(x)$ и получаете полную зависимость от времени. 

Раз вы знаете зависимость от времени, следовательно, вы знаете среднее значение любой динамической переменной. В квантовой механике динамической переменной соответствует эрмитов оператор $\hat{A}$, который действует в гильбертовом пространстве. Вы можете найти среднее значение для любой переменной, поместив оператор $\hat{A}$ между функциями состояния:
\begin{equation}
    \langle A \rangle = \langle \Psi | \hat{A} | \Psi \rangle
\end{equation}
Вы просто берете и считаете интеграл по всем координатам:
\begin{equation}
    \langle A \rangle = \int \Psi^*(x, t) \hat{A} \Psi(x, t) dx \label{eq:pure_mean}
\end{equation}
где интегрирование по $dx$ подразумевает интегрирование по всем координатам частиц. Функция $\Psi$ у вас нормирована на единицу, поэтому интеграл от квадрата модуля $\Psi$ равен единице. 

Теперь давайте рассмотрим множество тождественных копий нашей системы. Эти копии не взаимодействуют друг с другом. 
%[МЕСТО ДЛЯ ГРАФИКА: Таймкод ~ 29:45. Лектор рисует в ряд несколько одинаковых квадратных ящиков, под каждым из которых подписывает состояние $|\Psi\rangle$, и объединяет их фигурной скобкой, подписывая "чистый ансамбль"].
Каждая из этих копий находится в одном и том же состоянии $|\Psi\rangle$. Такая совокупность систем называется ансамблем. А так как все копии находятся в одном и том же квантовом состоянии, то это называется \textbf{чистый ансамбль}. 
% --- КОНЕЦ БЛОКА 2.2 ---

% --- НАЧАЛО БЛОКА 2.3 ---
\subsection{Проекционный оператор чистого состояния}

Оказывается, для чистого ансамбля удобно развить несколько другой математический аппарат. Мы вводим проекционный оператор $\hat{P}$, который определяется следующим образом:
\begin{equation}
    \hat{P} = |\Psi\rangle \langle \Psi|
\end{equation}
Давайте поясним, как этот оператор выглядит в матричном виде, если вы работаете в пространстве столбцов. В пространстве комплексных столбцов $\mathbb{E}^2$, где у вас есть базисные состояния $|+\rangle$ и $|-\rangle$, проекционный оператор $\hat{P}_1$ на состояние $|+\rangle$ задается матрицей:
\begin{equation}
    \hat{P}_1 = |+\rangle\langle+| = \begin{pmatrix} 1 & 0 \\ 0 & 0 \end{pmatrix}
\end{equation}
А проекционный оператор $\hat{P}_2$ на состояние $|-\rangle$ задается матрицей:
\begin{equation}
    \hat{P}_2 = |-\rangle\langle-| = \begin{pmatrix} 0 & 0 \\ 0 & 1 \end{pmatrix}
\end{equation}
Если вы подействуете оператором $\hat{P}_2$ на произвольный вектор-столбец $\begin{pmatrix} a \\ b \end{pmatrix}$, то матрица "вырежет" из этого состояния компоненту, соответствующую вектору $|-\rangle$, и вы получите вектор $\begin{pmatrix} 0 \\ b \end{pmatrix}$. Это и есть суть действия проекционных операторов.
% --- КОНЕЦ БЛОКА 2.3 ---

% --- НАЧАЛО БЛОКА 2.4 ---
\subsection{Свойства проекционного оператора}

Каковы основные свойства проекционного оператора $\hat{P}$?

Первое свойство: квадрат оператора равен самому оператору:
\begin{equation}
    \hat{P}^2 = \hat{P}
\end{equation}
Это логично: если вы еще раз подействуете оператором проекции на вектор, который уже спроектирован на данное состояние (и направлен вдоль него), вы получите тот же самый вектор.

Второе свойство: оператор является эрмитовым. Операция эрмитова сопряжения не меняет оператор:
\begin{equation}
    \hat{P}^\dagger = \hat{P}
\end{equation}

Третье, очень важное свойство: шпур (след) оператора $\hat{P}$ равен единице:
\begin{equation}
    Sp(\hat{P}) = 1
\end{equation}
Вы можете показать это напрямую. В общем случае шпур оператора — это сумма его диагональных матричных элементов. В координатном представлении с непрерывным индексом $x$ сумма заменяется на интеграл:
\begin{equation}
    Sp(\hat{P}) = \int dx \langle x | \hat{P} | x \rangle
\end{equation}
Подставим сюда определение оператора $\hat{P} = |\Psi\rangle \langle \Psi|$:
\begin{equation}
    Sp(\hat{P}) = \int dx \langle x | \Psi \rangle \langle \Psi | x \rangle
\end{equation}
Что такое $\langle x | \Psi \rangle$? Это вектор состояния $|\Psi\rangle$, спроектированный на базис координат $|x\rangle$. По определению, это и есть волновая функция $\Psi(x)$. А сомножитель $\langle \Psi | x \rangle$ — это комплексно сопряженная функция $\Psi^*(x)$. Поэтому мы получаем:
\begin{equation}
    Sp(\hat{P}) = \int \Psi(x) \Psi^*(x) dx = \int |\Psi(x)|^2 dx = 1
\end{equation}
Так как функция $\Psi$ нормирована на единицу, интеграл равен единице. Третье свойство доказано.
% --- КОНЕЦ БЛОКА 2.4 ---

% --- НАЧАЛО БЛОКА 2.5 ---
\subsection{Вычисление квантовомеханического среднего через шпур (след)}

Самое главное, четвертое свойство, отвечает на вопрос: для чего мы вообще вводили этот оператор? Оказывается, что среднее значение динамической переменной $\langle A \rangle$, заданное уравнением (\ref{eq:pure_mean}), можно эквивалентно записать через шпур произведения оператора наблюдаемой $\hat{A}$ и проекционного оператора $\hat{P}$:
\begin{equation}
    \langle A \rangle = \langle \Psi | \hat{A} | \Psi \rangle = Sp(\hat{A}\hat{P}) \label{eq:trace_mean}
\end{equation}
Эта форма записи очень важна, и давайте мы докажем, что это так и есть. Начнем с правой части — посчитаем шпур $Sp(\hat{A}\hat{P})$ в координатном $x$-представлении. Это снова интеграл от диагональных элементов:
\begin{equation}
    Sp(\hat{A}\hat{P}) = \int dx \langle x | \hat{A}\hat{P} | x \rangle
\end{equation}
Подставляем явный вид оператора $\hat{P}$:
\begin{equation}
    Sp(\hat{A}\hat{P}) = \int dx \langle x | \hat{A} | \Psi \rangle \langle \Psi | x \rangle
\end{equation}
Посмотрим внимательно на величину $\langle \Psi | x \rangle$. Это просто число (комплексная функция $\Psi^*(x)$). Я имею полное право переставить это число в начало подынтегрального выражения:
\begin{equation}
    Sp(\hat{A}\hat{P}) = \int dx \langle \Psi | x \rangle \langle x | \hat{A} | \Psi \rangle
\end{equation}
А теперь посмотрите, как мы эту сумму (интеграл) свернем. Из квантовой механики мы знаем условие полноты базиса. Единичный оператор $\hat{1}$ можно представить как интеграл по всем состояниям непрерывного координатного базиса:
\begin{equation}
    \hat{1} = \int dx | x \rangle \langle x |
\end{equation}
Это является обобщением на непрерывный случай известного соотношения для дискретного базиса $\sum_k |\Psi_k\rangle\langle\Psi_k| = \hat{1}$ (когда матрицы проекторов в сумме дают единичную матрицу). Подставляя это условие полноты $\int dx | x \rangle \langle x |$ в наше выражение, мы «схлопываем» базисные векторы $|x\rangle$:
\begin{equation}
    Sp(\hat{A}\hat{P}) = \langle \Psi | \left( \int dx | x \rangle \langle x | \right) \hat{A} | \Psi \rangle = \langle \Psi | \hat{1} \cdot \hat{A} | \Psi \rangle = \langle \Psi | \hat{A} | \Psi \rangle
\end{equation}
Всё тождественно совпало! Мы получили искомое среднее значение. Равенство доказано.

Поэтому форма вычисления среднего через волновую функцию $\langle \Psi | \hat{A} | \Psi \rangle$ оказывается полностью эквивалентна записи через шпур $Sp(\hat{A}\hat{P})$. Запомните это! Хотя вы больше привыкли к вычислению через волновые функции, форма со шпуром оказывается значительно удобнее. Именно она естественным образом обобщается на то, что мы сейчас введем для описания более сложных термодинамических ситуаций.
% --- КОНЕЦ БЛОКА 2.5 ---

% ==============================================================================
% РАЗДЕЛ 3
% ==============================================================================
\section{Смешанные ансамбли и равновесная матрица плотности}

А теперь давайте посмотрим на ситуацию шире. Все системы, макроскопические тела, вообще говоря, открыты. Мы все с чем-то взаимодействуем. Взаимодействие может быть слабеньким, может быть сильным, но система редко бывает полностью изолированной. И оказывается, что в таком случае система вообще не описывается какой-то одной волновой функцией. Как в этом случае описывать состояние квантовой системы?

% --- НАЧАЛО БЛОКА 3.1 ---
\subsection{Физические причины возникновения смешанных состояний}

%[МЕСТО ДЛЯ ГРАФИКА: Изображение макроскопической системы внутри резервуара. Лектор рисует замкнутую область "Система" и окружающую её бОльшую область "Резервуар". Таймкод: ~37:31]

Вот у нас имеется макроскопическая система, и она находится в резервуаре (термостате). Они по определению всегда взаимодействуют друг с другом. Раз взаимодействие есть, вы не можете сказать, что полная волновая функция составной системы «система + резервуар» факторизуется. Она не распадается в простое произведение волновой функции системы на волновую функцию резервуара:
\begin{equation}
    \Psi_{\text{сист+резервуар}} \neq \Psi_{\text{сист}} \cdot \Psi_{\text{резервуар}}
\end{equation}
Вы с этим явлением уже много раз встречались в квантовой механике. Какой самый разительный пример, когда подсистема не описывается отдельной волновой функцией? Вспомните систему, состоящую из протона и электрона (атом водорода). 

%[МЕСТО ДЛЯ ГРАФИКА: Рисунок атома водорода. Протон в центре, электрон на орбите. Таймкод: ~38:05]

Рассмотрим волновую функцию, соответствующую спиновому состоянию с полным спином $S = 0$ и проекцией спина $M_S = 0$ (синглетное состояние). Чему равна такая волновая функция системы «протон + электрон»? Она имеет вид антисимметричной суперпозиции:
\begin{equation}
    \Psi_{p+e} = \frac{1}{\sqrt{2}} \Big( |+\rangle_p |-\rangle_e - |-\rangle_p |+\rangle_e \Big)
\end{equation}
Можете ли вы в этом случае приписать какую-то отдельную волновую функцию только электрону? Нет, не можете. Вы можете приписать волновую функцию всей системе в целом, но про электрон вы не можете сказать, что он находится в каком-то конкретном квантовом состоянии (ни в $|+\rangle$, ни в $|-\rangle$). Отдельно взятый электрон здесь волновой функцией не описывается! Наша система, помещенная в резервуар, ведет себя точно так же: из-за квантовой запутанности с термостатом она не имеет собственной волновой функции.
% --- КОНЕЦ БЛОКА 3.1 ---

% --- НАЧАЛО БЛОКА 3.2 ---
\subsection{Смешанный ансамбль и двойное усреднение}

Что же делать в этом случае? Как описывать систему? Мы говорим следующее: если бы система была изолирована, у нее имелся бы набор стационарных квантовых состояний $|\Psi_k\rangle$. Мы не можем сказать, что система находится в каком-то одном из этих состояний или в их когерентной суперпозиции. Но мы можем сказать другое: \textit{определены вероятности} $w_k$ найти нашу систему в конкретном квантовом состоянии $|\Psi_k\rangle$.

Набор этих вероятностей $w_k$ определяет некоторое статистическое распределение. Чтобы это формализовать, мы вводим понятие \textbf{смешанного ансамбля}. 

%[МЕСТО ДЛЯ ГРАФИКА: Лектор рисует множество квадратиков (копий системы). Внутри каждого квадратика подписывает различные квантовые состояния: |\Psi_{i1}>, |\Psi_{i2}>, |\Psi_{i3}> и т.д. Таймкод: ~39:40]

Мы рассматриваем огромное множество мысленных копий нашей макроскопической системы. Каждая копия находится в каком-то одном из квантовых стационарных состояний. Первая — в $|\Psi_{i_1}\rangle$, вторая — в $|\Psi_{i_2}\rangle$ и так далее. Это множество копий и есть смешанный ансамбль. При этом доля систем в ансамбле, находящихся в состоянии $|\Psi_k\rangle$, в точности равна нашей вероятности $w_k$. Разумеется, сумма всех вероятностей равна единице:
\begin{equation}
    \sum_k w_k = 1
\end{equation}

Как в смешанном ансамбле вычислить среднее значение произвольной динамической переменной $A$? Нам придется сделать \textit{двойное усреднение}. Сначала мы берем квантовомеханическое среднее переменной $\hat{A}$ в заданном состоянии $|\Psi_k\rangle$, а затем усредняем эти значения по всем возможным состояниям ансамбля с весами $w_k$. То есть:
\begin{equation}
    \langle A \rangle = \sum_k w_k \langle \Psi_k | \hat{A} | \Psi_k \rangle
\end{equation}
Здесь $\langle \Psi_k | \hat{A} | \Psi_k \rangle$ — это чисто квантовое среднее, а сумма с весами $w_k$ — это уже статистическое усреднение.
% --- КОНЕЦ БЛОКА 3.2 ---

% --- НАЧАЛО БЛОКА 3.3 ---
\subsection{Определение статистического оператора (матрицы плотности)}

Оказывается, вместо формулы с двойным усреднением гораздо удобнее использовать несколько другую математическую запись. Для описания смешанного состояния вводится \textbf{статистический оператор}, синонимом которого является термин \textbf{матрица плотности}. Обозначается он буквой $\hat{\rho}$.

Определяется этот оператор следующим образом. Это взвешенная сумма проекционных операторов на каждое из стационарных состояний:
\begin{equation}
    \hat{\rho} = \sum_k w_k |\Psi_k\rangle \langle \Psi_k | \label{eq:density_matrix_def}
\end{equation}
В базисе собственных стационарных состояний исследуемой системы $|\Psi_{k}\rangle$ эта матрица имеет предельно простой вид. Если мы посчитаем матричный элемент $\langle \Psi_n | \hat{\rho} | \Psi_m \rangle$, мы увидим, что матрица диагональна, а на главной диагонали стоят как раз наши вероятности $w_k$:
\begin{equation}
    \hat{\rho} = 
    \begin{pmatrix}
        w_1 & 0 & 0 & \dots \\
        0 & w_2 & 0 & \dots \\
        0 & 0 & w_3 & \dots \\
        \vdots & \vdots & \vdots & \ddots
    \end{pmatrix}
\end{equation}
Матрица плотности — это инструмент, который заменяет собой волновую функцию в тех случаях, когда волновая функция неприменима.
% --- КОНЕЦ БЛОКА 3.3 ---

% --- НАЧАЛО БЛОКА 3.4 ---
\subsection{Свойства матрицы плотности}

Давайте перечислим основные свойства оператора $\hat{\rho}$. Вы увидите, что они очень похожи на свойства проекционного оператора для чистого состояния, но есть и фундаментальное отличие.

\textbf{1. Эрмитовость:} 
Оператор матрицы плотности эрмитов. 
\begin{equation}
    \hat{\rho}^\dagger = \hat{\rho}
\end{equation}
Это очевидно, так как проекционные операторы $|\Psi_k\rangle \langle \Psi_k|$ эрмитовы, а числа $w_k$ — это действительные вероятности.

\textbf{2. Нормировка (шпур равен единице):}
След (или шпур) оператора $\hat{\rho}$ равен сумме его диагональных элементов. А так как на диагонали стоят вероятности, сумма которых равна 1, получаем:
\begin{equation}
    Sp(\hat{\rho}) = \sum_k w_k = 1
\end{equation}

\textbf{3. Смысл диагональных элементов:}
Диагональные элементы матрицы плотности в любом представлении имеют смысл вероятностей (или плотности вероятности) обнаружить систему в соответствующем состоянии. 
\begin{equation}
    \langle \Psi_k | \hat{\rho} | \Psi_k \rangle = w_k
\end{equation}

\textbf{4. Отличие смешанного состояния от чистого:}
В общем случае для смешанного ансамбля квадрат оператора плотности \textit{не равен} самому оператору:
\begin{equation}
    \hat{\rho}^2 \neq \hat{\rho}
\end{equation}
Почему? Если мы возведем диагональную матрицу в квадрат, мы получим на диагонали вероятности в квадрате $w_k^2$. Поскольку вероятности $0 \le w_k \le 1$, то $w_k^2 < w_k$ (если $w_k \neq 1$). Равенство $\hat{\rho}^2 = \hat{\rho}$ выполняется \textit{только} тогда, когда одна вероятность равна $1$, а все остальные $0$. А это означает, что ансамбль сводится к единственному состоянию, то есть становится чистым. Таким образом, равенство $\hat{\rho}^2 = \hat{\rho}$ является строгим критерием того, что состояние системы чистое. 

\textbf{5. Вычисление статистического среднего:}
Самое главное свойство — как использовать этот аппарат. Оказывается, двойное усреднение полностью сводится к взятию шпура от произведения оператора наблюдаемой $\hat{A}$ и матрицы плотности $\hat{\rho}$:
\begin{equation}
    \langle A \rangle = Sp(\hat{A}\hat{\rho})
\end{equation}
Давайте докажем это равенство, начав с правой части. По определению, вычислим шпур в базисе состояний $|\Psi_{k'}\rangle$:
\begin{equation}
    Sp(\hat{A}\hat{\rho}) = \sum_{k'} \langle \Psi_{k'} | \hat{A} \hat{\rho} | \Psi_{k'} \rangle
\end{equation}
Подставим сюда определение матрицы плотности (\ref{eq:density_matrix_def}):
\begin{equation}
    Sp(\hat{A}\hat{\rho}) = \sum_{k'} \langle \Psi_{k'} | \hat{A} \left( \sum_k w_k |\Psi_k\rangle \langle \Psi_k | \right) | \Psi_{k'} \rangle
\end{equation}
Переставим суммы и вынесем числа $w_k$:
\begin{equation}
    Sp(\hat{A}\hat{\rho}) = \sum_{k'} \sum_k w_k \langle \Psi_{k'} | \hat{A} | \Psi_k \rangle \langle \Psi_k | \Psi_{k'} \rangle
\end{equation}
Поскольку наши состояния образуют ортонормированный базис, скалярное произведение $\langle \Psi_k | \Psi_{k'} \rangle$ — это символ Кронекера $\delta_{kk'}$. Он обнуляет все слагаемые в двойной сумме, кроме тех, где $k = k'$:
\begin{equation}
    Sp(\hat{A}\hat{\rho}) = \sum_k w_k \langle \Psi_k | \hat{A} | \Psi_k \rangle
\end{equation}
Но это выражение — не что иное, как наше изначальное определение двойного усреднения $\langle A \rangle$! Равенство доказано. Мы видим, что формула со шпуром работает и для смешанных ансамблей.
% --- КОНЕЦ БЛОКА 3.4 ---

% --- НАЧАЛО БЛОКА 3.5 ---
\subsection{Квантовое каноническое распределение}

Теперь посмотрим на применение этого аппарата. Мы уже обсуждали классическое распределение Гиббса и квантовое микроканоническое распределение. В микроканоническом распределении вероятности $w_k$ задаются просто: это $1/\Omega(E)$ для состояний внутри узкого энергетического слоя, и $0$ вне него. 

Нас же сейчас интересует более важный случай — \textbf{квантовое каноническое распределение} (система в термостате). Для канонического распределения вероятность найти систему в $k$-том квантовом состоянии дается формулой Гиббса:
\begin{equation}
    w_k = \frac{1}{Z} e^{-\frac{E_k}{T}}
\end{equation}
Подставим эти вероятности в наше общее определение матрицы плотности:
\begin{equation}
    \hat{\rho} = \sum_k w_k |\Psi_k\rangle \langle \Psi_k | = \sum_k \frac{1}{Z} e^{-\frac{E_k}{T}} |\Psi_k\rangle \langle \Psi_k |
\end{equation}
Давайте посмотрим на член $e^{-E_k/T} |\Psi_k\rangle$. Так как $|\Psi_k\rangle$ — это собственное состояние Гамильтониана $\hat{H}$ с собственным значением энергии $E_k$, то применение функции от Гамильтониана эквивалентно умножению на функцию от энергии:
\begin{equation}
    e^{-\frac{\hat{H}}{T}} |\Psi_k\rangle = e^{-\frac{E_k}{T}} |\Psi_k\rangle
\end{equation}
Мы можем воспользоваться этим и заменить число $e^{-E_k/T}$ на оператор $e^{-\hat{H}/T}$, действующий на вектор:
\begin{equation}
    \hat{\rho} = \frac{1}{Z} \sum_k e^{-\frac{\hat{H}}{T}} |\Psi_k\rangle \langle \Psi_k |
\end{equation}
Теперь мы можем вынести оператор $e^{-\hat{H}/T}$ за знак суммы, так как суммирование идет по состояниям:
\begin{equation}
    \hat{\rho} = \frac{1}{Z} e^{-\frac{\hat{H}}{T}} \left( \sum_k |\Psi_k\rangle \langle \Psi_k | \right)
\end{equation}
Но выражение в скобках $\sum_k |\Psi_k\rangle \langle \Psi_k |$ — это в точности разложение единичного оператора $\hat{1}$ (условие полноты базиса)! Умножение на единичный оператор ничего не меняет. В итоге мы получаем фундаментальное соотношение:
\begin{equation}
    \hat{\rho} = \frac{1}{Z} e^{-\frac{\hat{H}}{T}}
\end{equation}
Это \textbf{статистический оператор квантового канонического распределения}. Подчеркните его двумя чертами — это основа основ квантовой статистической физики.

Осталось только понять, как вычислять статсумму $Z$. Мы знаем, что шпур матрицы плотности обязан равняться единице: $Sp(\hat{\rho}) = 1$. Подставляя сюда наш оператор, получаем:
\begin{equation}
    Sp\left( \frac{1}{Z} e^{-\frac{\hat{H}}{T}} \right) = 1 \implies Z = Sp\left( e^{-\frac{\hat{H}}{T}} \right)
\end{equation}
Итак, статистическая сумма квантовой системы равна шпуру от экспоненты минус Гамильтониан делить на температуру. Именно отсюда мы в дальнейшем будем выводить большое каноническое распределение и строить всю квантовую термодинамику. На этом мы ставим точку на сегодня, в следующий раз продолжим развивать этот аппарат.
% --- КОНЕЦ БЛОКА 3.5 ---

\end{document}
